\documentclass[a4paper,14pt,oneside]{extarticle}

\usepackage[utf8]{inputenc}
\usepackage[T2A]{fontenc}
\usepackage[english,russian]{babel}
\usepackage[left=3cm,right=1.5cm,top=2cm,bottom=2cm]{geometry}
\usepackage{setspace}
\usepackage{graphicx}
\usepackage{float}
\usepackage{enumitem}
\usepackage{array}
\usepackage{longtable}
\usepackage[hidelinks]{hyperref}
\usepackage{caption}

\setstretch{1.35}
\setlength{\parindent}{1.25cm}
\setlength{\parskip}{0pt}
\addto\captionsrussian{\renewcommand{\contentsname}{Содержание}\renewcommand{\figurename}{Рисунок}}
\captionsetup{labelsep=endash,justification=centering}
\graphicspath{{./}}

\newcommand{\labimage}[2]{
\begin{figure}[H]
\centering
\includegraphics[width=\linewidth,height=0.68\textheight,keepaspectratio]{#1}
\caption{#2}
\end{figure}
}

\newcommand{\theoryplaceholder}[2]{
\begin{center}
\fbox{\begin{minipage}[c][4.2cm][c]{0.86\textwidth}
\centering
\textit{МЕСТО ДЛЯ СКРИНШОТА}\\
\small #1
\end{minipage}}

\vspace{0.15cm}
\textbf{Рисунок #2}
\end{center}
}

\begin{document}
\selectlanguage{russian}
\renewcommand{\contentsname}{Содержание}
\renewcommand{\figurename}{Рисунок}

\begin{titlepage}
\begin{center}
\textbf{МИНИСТЕРСТВО НАУКИ И ВЫСШЕГО ОБРАЗОВАНИЯ РФ}\\
\vspace{0.4cm}
\textbf{[Наименование образовательной организации]}\\
\textbf{[Наименование факультета / кафедры]}

\vfill

\textbf{\Large ОТЧЁТ}\\
\vspace{0.4cm}
\textbf{\large по лабораторной работе № 1}\\
\vspace{0.4cm}
\textbf{\large «Основы компьютерных сетей. Физический и канальный уровни модели OSI/ISO. Знакомство с Cisco Packet Tracer»}

\vfill

\begin{flushright}
\begin{minipage}{0.48\textwidth}
Выполнил(а): студент(ка) группы \underline{\hspace{3cm}}\\
\underline{\hspace{6cm}}\\
\vspace{0.3cm}
Проверил(а): \underline{\hspace{4cm}}
\end{minipage}
\end{flushright}

\vfill

\textbf{[Город]}\\
\textbf{2026}
\end{center}
\end{titlepage}

\tableofcontents
\newpage

\section{Цель и содержание работы}

Цель лабораторной работы --- изучить назначение физического и канального уровней
модели OSI/ISO, рассмотреть основные сетевые топологии и среды передачи данных,
освоить интерфейс Cisco Packet Tracer и базовые команды операционной системы
Cisco IOS.

В работе выполнены следующие практические блоки:
\begin{enumerate}
\item знакомство с интерфейсом Packet Tracer и представлением компьютерной сети;
\item изучение режимов работы и структуры команд командной строки Cisco IOS;
\item установка первоначальных параметров безопасности коммутатора;
\item просмотр ARP-кэша и анализ обмена ARP-сообщениями;
\item отработка комплексных практических навыков настройки устройств;
\item изучение физических характеристик и модулей межсетевых устройств.
\end{enumerate}

\section{Теоретические сведения}

\subsection{Модель OSI/ISO и стек TCP/IP}

Модель OSI/ISO описывает передачу данных как взаимодействие семи уровней:
физического, канального, сетевого, транспортного, сеансового, представления и
прикладного. Физический уровень передаёт поток битов по линии связи. Канальный
уровень формирует кадры, использует MAC-адресацию и контролирует целостность
передачи в пределах локальной сети. Сетевой уровень отвечает за логическую
адресацию и маршрутизацию.

\theoryplaceholder{Вставить схему семи уровней модели OSI/ISO.}{Т1 --- Модель OSI/ISO.}

Транспортный уровень обеспечивает доставку данных между процессами и использует
порты. На сеансовом уровне устанавливается и завершается сеанс взаимодействия.
Уровень представления отвечает за представление и преобразование данных, а
прикладной уровень предоставляет сетевые службы приложениям.

В стеке TCP/IP верхние три уровня OSI объединены в прикладной уровень, сетевой
уровень представлен уровнем Интернета, транспортный уровень сохраняется, а
физический и канальный уровни объединены в уровень доступа к сети. Передача
данных сопровождается инкапсуляцией: каждый нижележащий уровень добавляет к
данным управляющую информацию.

\theoryplaceholder{Вставить схему прохождения данных сверху вниз и формирования PDU.}{Т2 --- Инкапсуляция данных.}

\theoryplaceholder{Вставить сравнительную схему OSI/ISO и стека TCP/IP (модели DoD).}{Т3 --- OSI/ISO и TCP/IP.}

\subsection{Физический уровень и топологии}

Физический уровень определяет среду передачи, электрические или оптические
сигналы, скорость, способ кодирования, разъёмы и распиновку. В локальных сетях
используются витая пара, коаксиальный и оптоволоконный кабели, а также
беспроводная среда. Для Fast Ethernet обычно достаточно двух пар витой пары,
для Gigabit Ethernet используются все четыре пары.

К основным физическим топологиям относятся шина, звезда, кольцо, ячеистая и
смешанная топологии. В современных Ethernet-сетях наиболее распространена
топология звезда, в центре которой находится коммутатор. Отказ одного кабеля
обычно изолирует только подключённое устройство, а отказ центрального
коммутатора нарушает работу сегмента.

\theoryplaceholder{Вставить схему выбранной физической топологии сети или сравнение топологий.}{Т4 --- Физические топологии сети.}

\subsection{Канальный уровень, Ethernet и MAC-адресация}

Канальный уровень преобразует поток битов в кадры, выполняет адресацию,
обнаружение ошибок и управление доступом к среде передачи. В стандартах IEEE
802 он разделён на подуровни LLC и MAC. Подуровень LLC связывает сетевой и
канальный уровни, а MAC отвечает за инкапсуляцию данных и доступ к среде.

MAC-адрес Ethernet имеет длину 48 бит и записывается в виде двенадцати
шестнадцатеричных цифр. Первые 24 бита содержат идентификатор производителя
OUI, последние 24 бита задаются производителем. Используются одноадресные,
широковещательные и многоадресные MAC-адреса. Широковещательный адрес Ethernet
имеет вид \texttt{FF-FF-FF-FF-FF-FF}.

Коммутатор работает преимущественно на втором уровне OSI. Он изучает MAC-адрес
источника каждого кадра и строит таблицу соответствия MAC-адресов портам.
Известные одноадресные кадры передаются только на нужный порт, а кадры с
неизвестным адресом назначения и широковещательные кадры распространяются по
другим портам сегмента.

\theoryplaceholder{Вставить изображение распиновки кабеля RJ-45 и обозначения пар.}{Т5 --- Распиновка RJ-45.}

\theoryplaceholder{Вставить схему Ethernet с физическим и канальным уровнями.}{Т6 --- Ethernet.}

\theoryplaceholder{Вставить перечень или схему основных стандартов IEEE Ethernet.}{Т7 --- Стандарты IEEE Ethernet.}

\theoryplaceholder{Вставить схему подуровня MAC Ethernet и его функций.}{Т8 --- Подуровень MAC Ethernet.}

\theoryplaceholder{Вставить схему работы CSMA/CD и обнаружения коллизии.}{Т9 --- Ассоциативный доступ и CSMA/CD.}

\theoryplaceholder{Вставить схему структуры Ethernet MAC-адреса: OUI и серийная часть.}{Т10 --- Структура MAC-адреса Ethernet.}

\subsection{Cisco IOS и ARP}

Cisco IOS предоставляет иерархический интерфейс командной строки. Пользовательский
режим обозначается приглашением \texttt{Switch>}, привилегированный --- 
\texttt{Switch\#}, глобальный режим конфигурации --- \texttt{Switch(config)\#}.
Команда \texttt{enable} переводит устройство в привилегированный режим, а
\texttt{configure terminal} открывает глобальный режим конфигурации. Для контроля
состояния используются команды семейства \texttt{show}, например
\texttt{show running-config}, \texttt{show startup-config} и
\texttt{show interfaces}.

Протокол ARP сопоставляет IPv4-адрес с MAC-адресом в локальном сегменте.
Результаты сопоставления сохраняются в ARP-кэше. В Windows содержимое кэша
просматривается командой \texttt{arp -a}, отдельные записи удаляются командой
\texttt{arp -d <адрес>}, а весь кэш очищается командой \texttt{arp -d *}.

\theoryplaceholder{Вставить пример физического MAC-адреса сетевого адаптера и форматы его записи.}{Т11 --- Физический MAC-адрес узла.}

\theoryplaceholder{Вставить схему одноадресной передачи: IP- и MAC-адреса источника и назначения.}{Т12 --- Одноадресная рассылка.}

\theoryplaceholder{Вставить схему широковещательной передачи и адрес FF-FF-FF-FF-FF-FF.}{Т13 --- Широковещательная рассылка.}

\theoryplaceholder{Вставить схему многоадресной передачи и соответствия IPv4-группы MAC-адресу.}{Т14 --- Многоадресная рассылка.}

В интерфейсе IOS контекстная справка вызывается знаком вопроса
\texttt{?}. Она показывает допустимые команды и параметры с учётом текущего
режима и уже введённой части команды. Клавиши Tab, стрелки истории,
Ctrl-Z и Ctrl-C ускоряют настройку и навигацию по CLI.

\theoryplaceholder{Вставить скриншот контекстной справки IOS после команды clock set.}{Т15 --- Контекстная справка Cisco IOS.}

\section{Ход выполнения и результаты}

Представленные скриншоты содержат как условия практических упражнений, так и
рабочие области Packet Tracer с индикаторами состояния и таблицами
оценивания. Поэтому ниже для каждого изображения указано, какое действие
требуется выполнить или что именно проверяется на экране. Для скриншотов с
таблицами оценивания отдельно учитывается отображаемый прогресс задания.

\subsection{Знакомство с интерфейсом Cisco Packet Tracer}

В Cisco Packet Tracer были рассмотрены рабочая область, панели выбора сетевых
устройств, режимы реального времени и моделирования, а также инструменты
создания и исследования сети. В режиме моделирования можно проследить
прохождение PDU между устройствами и увидеть, какие уровни модели OSI
задействуются при обмене данными. Результаты выполнения задания представлены
на рисунках 1--3.

\subsection{Навигация по Cisco IOS}

Выполнено подключение к коммутатору через консольный кабель и изучены режимы
CLI. Проверены переходы командами \texttt{enable}, \texttt{configure terminal},
\texttt{exit}, \texttt{end} и \texttt{disable}. С помощью знака вопроса
\texttt{?} и клавиши Tab рассмотрены контекстная справка и автодополнение
команд. Команда \texttt{clock set} использовалась для изучения синтаксической
подсказки и установки системного времени. Результаты приведены на рисунках
4--7.

\subsection{Первоначальная настройка безопасности коммутатора}

Для коммутатора были выполнены базовые настройки:
\begin{itemize}
\item задано имя устройства;
\item настроен пароль консольной линии и включён запрос пароля;
\item задан пароль привилегированного режима;
\item настроен зашифрованный пароль \texttt{enable secret};
\item включено шифрование открытых паролей командой
\texttt{service password-encryption};
\item задано предупреждающее MOTD-сообщение;
\item конфигурация сохранена в NVRAM командой
\texttt{copy running-config startup-config}.
\end{itemize}

После настройки проверены вход в консоль, переход в привилегированный режим,
вывод текущей и сохранённой конфигурации. Результаты представлены на рисунках
8--12.

\subsection{Работа с ARP и анализ пакетов}

В рамках задания исследовалась таблица ARP операционной системы Windows.
Сначала ARP-кэш очищался, затем командой \texttt{ping} формировался запрос к
узлу или шлюзу. После этого команда \texttt{arp -a} позволяла убедиться в
появлении соответствующей записи.

В Wireshark для анализа применялся фильтр \texttt{arp}. Первый пакет обмена
является ARP-запросом широковещательного типа: отправитель знает IP-адрес
назначения, но ещё не знает его MAC-адрес. Второй пакет является ARP-ответом,
который сообщает MAC-адрес владельца запрошенного IP-адреса. После получения
ответа запись сохраняется в кэше, поэтому последующие пакеты к этому узлу не
требуют нового ARP-запроса до истечения времени хранения записи.

ARP добавляет небольшую задержку перед первым обменом, но кэширование уменьшает
число повторных запросов и нагрузку на локальную сеть. Иллюстрации выполненной
работы представлены на рисунках 13--14.

\subsection{Комплексная настройка сети}

В комплексном задании настроены имена и IPv4-адреса интерфейсов двух
коммутаторов, установлены пароли, добавлено MOTD-сообщение, сохранены
конфигурации и проверена связность между устройствами. Использовалась
адресация из сети \texttt{10.10.10.0/24}: для Room-145 применён адрес
\texttt{10.10.10.100}, для Room-146 --- \texttt{10.10.10.150}, для Manager ---
\texttt{10.10.10.4}, для Reception --- \texttt{10.10.10.5}. Топология и
таблица адресации показаны на рисунках 15--16.

\subsection{Изучение межсетевых устройств}

Изучены физические интерфейсы маршрутизатора, доступные слоты расширения и
совместимые модули. Для организации соединений выбирались соответствующие
медные прямые, перекрёстные и последовательные кабели. Также рассмотрено
подключение маршрутизаторов, коммутаторов и конечных устройств с учётом типа
порта и назначения интерфейса. Результаты приведены на рисунках 17--18.

\section{Пошаговое выполнение задания по методичке}

Ниже пункты приведены в той же последовательности, что и в методических
указаниях. Для каждого шага указаны действие, используемый инструмент и
проверяемый результат. Скриншоты содержат как условия упражнений, так и
рабочие области Packet Tracer с таблицами оценивания, поэтому фактические
значения, отсутствующие на изображениях, не подменяются предположениями.

\subsection{Задание 3.1. Знакомство с интерфейсом Packet Tracer}

\subsubsection{Часть 1. Обзор программы Packet Tracer}

\paragraph{Шаг 1. Получение доступа к материалам.}
Открыт Cisco Packet Tracer и просмотрены учебные материалы, доступные через
разделы справки и обучения. Определены рабочая область, панели выбора
устройств, область сценария и панель проверки результатов. Содержание
исходного упражнения показано на рисунке 1.

\paragraph{Шаг 2. Режимы работы Packet Tracer.}
Проверено переключение между режимами \textit{Realtime} и
\textit{Simulation}. В режиме реального времени сеть работает непрерывно,
а в режиме моделирования можно пошагово просматривать PDU, заголовки
протоколов и путь прохождения сообщения. Для этого используются фильтры
событий, \textit{Capture/Forward} и \textit{Auto Capture/Play}.

\paragraph{а. Рабочая область Realtime.}
Открыта рабочая область \textit{Realtime}. В этом режиме сеть работает
непрерывно, а изменения конфигурации применяются сразу.

\paragraph{б. Переход в Simulation.}
Открыта вкладка \textit{Simulation}; скорость работы сети уменьшилась,
что позволило наблюдать прохождение отдельных пакетов.

\paragraph{в. Auto Capture/Play.}
Запущено автоматическое моделирование кнопкой
\textit{Auto Capture/Play}. В рабочей области появились конверты PDU,
движущиеся между сетевыми устройствами.

\paragraph{г. Приостановка моделирования.}
Повторным нажатием \textit{Auto Capture/Play} моделирование остановлено,
чтобы изучить состояние пакетов на текущем шаге.

\paragraph{д. Пошаговый переход.}
Кнопкой \textit{Capture/Forward} выполнен переход на один шаг. Повторными
нажатиями прослеживалось прохождение сообщения через промежуточное
устройство.

\paragraph{е. Анализ PDU.}
Открыты сведения входящего и исходящего PDU и вкладка OSI Model.
Сравнено количество отображаемых уровней и содержание заголовков до и после
обработки пакета.

\paragraph{Шаг 3. Логический и физический виды.}
Проверены логическое представление топологии и физическое представление
расположения устройств. Логический вид показывает связи и конфигурацию
сети, физический --- комнаты, стойки, оборудование и кабельные соединения.

\paragraph{а. Logical.}
Открыта рабочая область \textit{Logical}, используемая для построения
топологии, настройки устройств и устранения неполадок.

\paragraph{б. Physical.}
Открыта рабочая область \textit{Physical}, в которой оцениваются размеры
и размещение компонентов сети, а также физическая топология.

\paragraph{в. Назначение физической области.}
Проверено, что физическая область помогает оценить размещение устройств
и кабельных соединений в реальной среде, но не заменяет логическую схему.

\paragraph{г. Возврат в Logical.}
Кнопкой переключения выполнен возврат из \textit{Physical} в рабочую область
\textit{Logical}.

\subsubsection{Часть 2. Локальные сети, WAN и Интернет}

\paragraph{Шаг 1. Компоненты сети.}
Рассмотрены категории \textit{End Devices}, \textit{Network Devices} и
\textit{Connections}. Оконечные устройства формируют или принимают данные,
промежуточные устройства пересылают трафик, а соединения задают физическую
среду передачи.

\paragraph{а. Категории промежуточных устройств.}
Перечислены категории промежуточных устройств: маршрутизаторы,
коммутаторы, концентраторы, точки доступа и модемы. В Packet Tracer
категория выбирается на панели сетевых устройств.

\paragraph{б. Перечень промежуточных устройств.}
В топологии отдельно отмечены маршрутизаторы, коммутаторы и устройства
доступа, соединяющие оконечные узлы или сетевые сегменты.

\paragraph{в. Оконечные устройства топологии.}
Без учёта облаков Internet и Intranet подсчитываются значки оконечных
устройств, причём каждый физический узел считается один раз, даже если у
него несколько подключений. К оконечным узлам относятся компьютеры,
серверы, принтеры и телефоны.

\paragraph{г. Промежуточные устройства топологии.}
Без учёта двух облаков подсчитываются значки устройств, имеющих несколько
входящих подключений. К ним относятся коммутаторы, маршрутизаторы и
другие устройства, пересылающие трафик.

\paragraph{д. Маршрутизаторы.}
Отдельно подсчитываются маршрутизаторы. Устройство Linksys в данном
упражнении также считается маршрутизатором, даже если оно отображается
в отдельной категории.

\paragraph{е. Оконечные устройства, не являющиеся ПК.}
Отдельно выделяются серверы, принтеры, IP-телефоны и другие оконечные
устройства, которые не являются настольными компьютерами.

\paragraph{ж. Типы соединений.}
Подсчитываются различные типы физических соединений, представленные в
топологии: медные, оптические, последовательные, консольные и
беспроводные, если они используются.

\paragraph{з. Беспроводная связь.}
В категории \textit{Connections} может отсутствовать отдельный значок
беспроводного кабеля, поскольку беспроводное соединение не является
физическим кабелем и настраивается через беспроводной интерфейс устройства.

\paragraph{Шаг 2а. Роль Server-PT.}
Server-PT может выступать в роли сервера, потому что предоставляет
сетевые службы другим узлам. Настольный или переносной компьютер также
может выполнять серверную программу, но в учебной модели его основная
роль обозначена как клиентская.

\paragraph{Шаг 2б. Модели клиент--сервер.}
Модель клиент--сервер означает, что клиент отправляет запрос, а сервер
обрабатывает его и возвращает результат. Примеры промежуточных устройств:
коммутатор и маршрутизатор.

\paragraph{Шаг 2в. Выбор среды передачи.}
Критериями выбора являются требуемая скорость, максимальная длина линии,
стоимость, устойчивость к помехам, тип интерфейса и необходимость
беспроводного доступа.

\paragraph{Шаг 3а. Отличия LAN и WAN.}
LAN объединяет устройства в ограниченной области, например в аудитории
или здании. WAN соединяет удалённые LAN через каналы оператора связи.

\paragraph{Шаг 3б. Глобальные сети в Packet Tracer.}
В топологии представлены глобальные облака Internet и Intranet, через
которые условно показываются внешние и корпоративные сети.

\paragraph{Шаг 3в. Локальные сети в Packet Tracer.}
Локальные сегменты представлены группами оконечных устройств,
подключённых к коммутаторам и маршрутизаторам внутри отдельных площадок.

\paragraph{Шаг 3г. Упрощение Интернета.}
Модель Интернета в Packet Tracer значительно упрощена: облако заменяет
множество реальных операторов, маршрутизаторов, автономных систем и
магистральных каналов.

\paragraph{Шаг 3д. Подключение домашних пользователей.}
Распространённые способы доступа: оптический канал, DSL, кабельный
модем, мобильная сеть и домашний Wi-Fi через маршрутизатор.

\paragraph{Шаг 3е. Подключение предприятий.}
Для предприятий применяются выделенные оптические линии, Ethernet
оператора, MPLS, VPN, радиорелейные и резервные WAN-каналы.

\paragraph{Сложные задачи части 2.}
\begin{enumerate}[label=\arabic*.]
\item Добавлено оконечное устройство в одну из локальных сетей и выбрана
среда передачи. Обосновано, что устройству нужен сетевой интерфейс,
например Ethernet или Wi-Fi. Правильность подключения проверяется по
активному индикатору порта и прохождению PDU.
\item Добавлено промежуточное устройство и подключено к локальной или
глобальной сети физическим носителем. Для него определено назначение:
коммутатор соединяет узлы LAN, а маршрутизатор соединяет разные сети.
\item Создан новый проект с двумя локальными сетями, соединёнными через
глобальную сеть. Все устройства подключены, схема проверена, ответы
сохранены в Packet Tracer. Таблица распределения баллов находится на
рисунке 3.
\end{enumerate}

\subsection{Задание 3.2. Режимы работы и команды CLI}

\subsubsection{Часть 1. Базовое подключение и справка}

\paragraph{Шаг 1. Подключение PC к коммутатору.}
К компьютеру подключён консольный кабель, выбран последовательный интерфейс
и открыт терминал Packet Tracer. Такое подключение используется для
первоначальной настройки устройства без IP-адреса.

\paragraph{а. Выбор категории Connections.}
В нижней части окна Packet Tracer открыта категория
\textit{Connections}, обозначенная значком молнии.

\paragraph{б. Выбор консольного кабеля.}
Выбран светло-голубой консольный кабель, предназначенный для подключения
компьютера к консольному порту сетевого устройства.

\paragraph{в. Выбор PC1 и RS-232.}
На PC1 выбран последовательный интерфейс RS-232 как сторона компьютера
для консольного подключения.

\paragraph{г. Выбор коммутатора S1.}
На коммутаторе S1 открыт список портов и выбран консольный порт, а не
Ethernet-порт передачи данных.

\paragraph{д. Завершение подключения.}
После выбора консольного порта соединение PC1--S1 появилось в рабочей
области, что подтверждает правильность выбора кабеля и интерфейсов.

\paragraph{Шаг 2. Запуск терминала.}
В окне терминала оставлены параметры последовательного подключения по
умолчанию и запущен сеанс. После этого появилась командная строка
коммутатора.

\paragraph{а. Открытие Desktop.}
На PC1 открыта вкладка \textit{Desktop}, в которой доступны приложения
для настройки и проверки сетевого подключения.

\paragraph{б. Запуск Terminal.}
Открыто приложение \textit{Terminal}, проверены параметры последовательного
подключения по умолчанию и подтверждён запуск.

\paragraph{в. Подтверждение параметров.}
Нажата кнопка \textit{OK}; после этого установлен консольный сеанс с S1.

\paragraph{г. Сообщение начала работы.}
В окне терминала появилось сообщение \textit{Press RETURN to get started!}.
После нажатия Enter отображается приглашение пользовательского режима
\texttt{S1>}.

\paragraph{Шаг 3. Контекстная справка.}
Проверено, что \texttt{?} выводит доступные команды текущего режима,
\texttt{s?} показывает команды, начинающиеся с введённых символов, а пробел
перед \texttt{?} показывает допустимые параметры уже набранной команды.
Клавиша Tab завершает однозначную команду. Последовательность показана на
рисунке 4.

\paragraph{а. Справка в пользовательском режиме.}
В приглашении \texttt{S1>} введён знак \texttt{?}; получен список команд,
доступных в пользовательском EXEC-режиме. По записи \texttt{S1>?} также
проверено, какие команды начинаются с буквы \texttt{s}.

\paragraph{б. Справка после команды.}
Введено начало команды с пробелом перед \texttt{?}. IOS показала допустимые
параметры и продолжения команды в текущем контексте.

\paragraph{в. Контекстная справка.}
Введён знак \texttt{?} после неполной команды без завершающего пробела.
Получен список команд и ключевых слов, начинающихся с указанной
последовательности символов.
\subsubsection{Часть 2. Изучение режимов ввода}

\paragraph{Шаг 1. Привилегированный режим.}
Командой \texttt{enable} выполнен переход из \texttt{Switch>} в
\texttt{Switch\#}. Команда \texttt{enable ?} использована для проверки
справки, а клавиша Tab --- для автодополнения.

\paragraph{а. Справка по enable.}
Командой \texttt{enable ?} проверен список допустимых параметров. Для
обычного перехода параметр не требуется.

\paragraph{б. Автодополнение Tab.}
После ввода \texttt{en} и нажатия Tab команда автоматически завершена до
\texttt{enable}. Это является завершением команды клавишей Tab.

\paragraph{в. Проверка приглашения.}
После выполнения \texttt{enable} приглашение изменилось на
\texttt{S1\#}, что подтверждает переход в привилегированный режим.

\paragraph{г. Поиск команд по первой букве.}
Введено \texttt{s?}; IOS вывела команды, начинающиеся с буквы \texttt{s}.
Это позволяет уточнять имя команды без полного набора.

\paragraph{Шаг 2. Глобальная конфигурация.}
Командой \texttt{configure terminal} выполнен переход в
\texttt{Switch(config)\#}, командой \texttt{exit} --- возврат. Изменение
приглашения служит проверкой текущего режима. Действия показаны на рисунке 5.

\paragraph{а. Команда configure.}
В привилегированном режиме введена команда \texttt{configure} или её
однозначное сокращение. IOS вывела запрос выбора источника конфигурации
\texttt{[terminal]}.

\paragraph{б. Подтверждение terminal.}
Нажата клавиша Enter для принятия параметра \texttt{terminal} по умолчанию.
Приглашение изменилось на \texttt{S1(config)\#}.

\paragraph{в. Возврат из режима конфигурации.}
Командой \texttt{exit}, \texttt{end} или сочетанием Ctrl-Z выполнен возврат
в привилегированный режим. Результат подтверждается приглашением
\texttt{S1\#}.

\subsubsection{Часть 3. Настройка часов}

\paragraph{Шаг 1. Справка и синтаксис.}
Вызваны \texttt{clock ?} и \texttt{clock set ?}; по подсказке определены
обязательные параметры времени, дня, месяца и года.

\paragraph{а. Просмотр часов.}
Командой \texttt{show clock} просмотрено текущее время и дата
коммутатора.

\paragraph{б. Неполная команда clock.}
Введена команда \texttt{clock}; IOS сообщила о необходимости дополнительных
параметров.

\paragraph{в. Контекстная справка clock.}
Командой \texttt{clock ?} получен список продолжений команды и подтверждено,
что IOS позволяет уточнять синтаксис непосредственно в CLI.

\paragraph{г. Команда clock set.}
Командой \texttt{clock set ?} последовательно определены параметры,
необходимые для установки времени.

\paragraph{д. Установка времени 15:00.}
Введено время \texttt{15:00:00} в 24-часовом формате и проверено наличие
дополнительных параметров даты.

\paragraph{е. Установка даты 31 января 2035 года.}
По подсказке введены день, месяц и год; затем командой \texttt{show clock}
проверена строка с установленными значениями.

\paragraph{ж. Исправление неверных данных.}
Если исходные данные отличаются от требуемых, повторно выполнена команда
\texttt{clock set 15:00:00 31 Jan 2035}, как указано в задании.

\paragraph{Шаг 2. Ошибки синтаксиса.}
Проверено, что неполная команда вызывает сообщение о недостающем параметре,
а неизвестная команда --- сообщение об ошибке. Проверены варианты
\texttt{clock}, \texttt{clock set 25:00:00} и команда с отсутствующим
параметром дня. Результат показан на рисунке 6.

\paragraph{а. Неполная команда.}
На запрос \texttt{clock} IOS возвращает сообщение
\textit{Incomplete command}, поскольку не указаны параметры.

\paragraph{б. Некорректное время.}
Команда \texttt{clock set 25:00:00} отклоняется как содержащая
недопустимое значение часа.

\paragraph{в. Недостающий параметр даты.}
Команда с временем и неполной датой также отклоняется до ввода всех
обязательных параметров.

\paragraph{г. Корректная команда.}
Полная команда с допустимыми временем, днём, месяцем и годом принимается,
после чего \texttt{show clock} показывает результат.

\paragraph{Шаг 2. Дополнительные сообщения.}
Проверены сообщения IOS при сокращённых и неверных командах. Таблица
оценивания частей задания приведена на рисунке 7.

\subsection{Задание 3.3. Первоначальная настройка безопасности}

\subsubsection{Часть 1. Конфигурация по умолчанию}

\paragraph{Шаг 1.}
Выполнен переход в привилегированный режим командой \texttt{enable}.

\paragraph{Шаг 2.}
Командой \texttt{show running-config} просмотрена текущая конфигурация.
Проверяются количество FastEthernet- и GigabitEthernet-интерфейсов, значения
параметров по умолчанию и настройки линий VTY.

\paragraph{а. FastEthernet-интерфейсы.}
В выводе \texttt{show running-config} подсчитано количество портов
FastEthernet. Конкретное число зависит от выбранной модели коммутатора и
переписывается из вывода Packet Tracer.

\paragraph{б. GigabitEthernet-интерфейсы.}
Аналогично подсчитано количество портов GigabitEthernet. В отчёте оно
определяется моделью устройства, показанной в рабочем окне.

\paragraph{в. Параметры VTY-линий.}
Проверены значения по умолчанию для виртуальных терминальных линий VTY,
через которые выполняется удалённый доступ к CLI.

\paragraph{г. Просмотр NVRAM.}
Текущую конфигурацию показывает \texttt{show running-config}, а содержимое
NVRAM --- \texttt{show startup-config}.

\paragraph{д. Отсутствие startup-config.}
Сообщение \texttt{startup-config is not present} появляется до первого
сохранения, потому что конфигурация ещё не записана в NVRAM. Исходная
проверка показана на рисунке 8.

\subsubsection{Часть 2. Основные параметры коммутатора}

\paragraph{а. Имя устройства.}
Задано имя устройства командой \texttt{hostname S1}.

\paragraph{б. Защита консоли.}
В режиме \texttt{line console 0} задан пароль консоли и включён
\texttt{login}.

\paragraph{в. Проверка консоли.}
Выполнено повторное подключение и проверен запрос
\textit{User Access Verification}.

\paragraph{г. Пароль enable.}
Задан пароль привилегированного режима
\texttt{enable password cisco}.

\paragraph{д. Проверка enable.}
Повторным переходом через \texttt{enable} проверен запрос пароля.

\paragraph{е. Защищённый пароль.}
Задан защищённый пароль \texttt{enable secret itsasecret}.

\paragraph{ж. Проверка running-config.}
Командой \texttt{show running-config} проверены внесённые настройки.

\paragraph{з. Шифрование паролей.}
Командой \texttt{service password-encryption} включено шифрование
открытых паролей.

Последовательность настройки консоли, паролей и MOTD показана на рисунках 9
и 10. Команда \texttt{enable secret} имеет приоритет над
\texttt{enable password}; команда шифрования дополнительно защищает пароли
линий и обычный пароль enable.

\subsubsection{Часть 3. Баннер MOTD}

Введена команда \texttt{banner motd "This is a secure system. Authorized
Access Only"}. При последующем подключении сообщение появляется перед
запросом пароля. Баннер предупреждает о разрешённом доступе и информирует
пользователя о правилах работы с устройством.

\subsubsection{Часть 4. Сохранение конфигурации}

\paragraph{а. Проверка конфигурации.}
Командой \texttt{show run} просмотрена текущая конфигурация.

\paragraph{б. Сохранение.}
Командой \texttt{copy running-config startup-config} выполнено
сохранение в NVRAM. Краткая форма --- \texttt{copy run start}.

\paragraph{в. Проверка startup-config.}
Командой \texttt{show startup-config} проверено наличие сохранённой
конфигурации.

\subsubsection{Часть 5. Настройка S2}

\paragraph{а. Имя устройства.}
Для второго коммутатора задано имя \texttt{S2}.

\paragraph{б. Защита консоли.}
Для \texttt{line console 0} задан пароль \texttt{letmein} и включён
\texttt{login}.

\paragraph{в. Пароли режимов.}
Для привилегированного режима задан пароль \texttt{c1\$c0}, а в качестве
\texttt{enable secret} задан пароль \texttt{itsasecret}.

\paragraph{г. Баннер MOTD.}
Настроено сообщение: \texttt{Authorized access only. Unauthorized access is
prohibited and violators will be prosecuted to the full extent of the law.}

\paragraph{д. Шифрование.}
Командой \texttt{service password-encryption} зашифрованы открытые пароли
консольной линии и обычный пароль enable.

\paragraph{е. Проверка конфигурации.}
Командой \texttt{show running-config} проверены имя, линии доступа, пароли,
баннер и сервис шифрования.

\paragraph{ж. Сохранение.}
Конфигурация S2 сохранена в NVRAM командой
\texttt{copy running-config startup-config}. Требования и порядок действий
отображены на рисунке 11, таблица оценивания --- на рисунке 12.

\subsection{Задание 3.4. ARP и Wireshark}

\subsubsection{Часть 1. Просмотр ARP-кэша}

\begin{enumerate}[label=\arabic*.]
\item Команда \texttt{arp -a} отображает все записи ARP-кэша.
\item Команда \texttt{arp -d *} удаляет все записи, а
\texttt{arp -d 192.168.1.11} удаляет запись только указанного адреса.
\item После \texttt{ping 192.168.1.2} запись добавляется динамически.
\item Физический адрес узла определяется повторной командой
\texttt{arp -a}. Конкретный MAC-адрес на приложенных скриншотах не показан,
поэтому он должен быть перенесён из фактического вывода Windows.
\end{enumerate}

\subsubsection{Часть 2. Ручная работа с кэшем}

\paragraph{а. Подготовка адресов.}
Удалены старые записи и отправлены ping на все адреса другого учащегося и
на адрес шлюза по умолчанию.

\paragraph{б. Проверка записей.}
Командой \texttt{arp -a} проверено, что IP-адреса, к которым обращались,
появились в кэше.

\paragraph{в. Повторный ping.}
Для отсутствующей записи повторно отправлен ping на адрес назначения, после
чего таблица ARP проверена ещё раз.

\paragraph{г. Запуск с правами администратора.}
Открыта командная строка с правами администратора, необходимыми для
удаления всех записей кэша.

\paragraph{д. Очистка всех записей.}
Выполнена команда \texttt{arp -d *}.

\paragraph{е. Контроль очистки.}
Командой \texttt{arp -a} проверено, что записи удалены или находятся в
состоянии повторного формирования.

\paragraph{ж. Повторное заполнение.}
После ожидания проверено появление динамических записей, создаваемых при
обращении к сетевым узлам.

\subsubsection{Часть 3. Захват сообщений ARP}

\paragraph{а. Выбор интерфейса.}
В Wireshark выбран активный интерфейс и начат захват пакетов. Установлен
фильтр отображения \texttt{arp}.

\paragraph{б. Очистка кэша.}
Перед захватом выполнена команда \texttt{arp -d *}, затем командой
\texttt{arp -a} проверено состояние кэша.

\paragraph{в. Ping шлюза.}
Выполнена команда \texttt{ping 192.168.1.1}.

\paragraph{г. Остановка захвата.}
После появления ответов ICMP захват остановлен и в списке оставлены
пакеты ARP.

\paragraph{д. Первый пакет.}
Первым пакетом является ARP-запрос с широковещательным адресом назначения.

\paragraph{е. Второй пакет.}
Вторым является ARP-ответ с MAC-адресом шлюза. Для обоих пакетов в отчёте
фиксируются MAC отправителя, IP отправителя, MAC назначения и IP назначения.
Числовые значения этих полей отсутствуют на приложенных скриншотах и должны
быть переписаны из панели сведений Wireshark.

\paragraph{ж. Фильтр ARP и ICMP.}
Для анализа задержки применён фильтр \texttt{arp or icmp}.

\paragraph{з. Повторный ping.}
После очистки кэша выполнена команда \texttt{ping 192.168.1.1}, затем
и компьютера другого учащегося.

\paragraph{и. Сравнение времени.}
Первый ping включает время ARP-разрешения, последующие используют запись
кэша и обычно передаются без дополнительного ARP-запроса.

\subsection{Задание 3.5. Skills Integration Challenge}

По таблице адресации и требованиям на рисунке 13 выполнены следующие пункты:
\begin{enumerate}[label=\arabic*.]
\item заданы имена двух коммутаторов;
\item назначены адреса VLAN-интерфейсам Room-145 и Room-146;
\item назначены адреса Manager и Reception;
\item установлены пароли консоли и привилегированного режима;
\item добавлено MOTD-сообщение;
\item включено шифрование паролей;
\item сохранены конфигурации в NVRAM;
\item проверена связность между всеми устройствами.
\end{enumerate}

В таблице использованы Room-145 --- \texttt{10.10.10.100/24}, Room-146 ---
\texttt{10.10.10.150/24}, Manager --- \texttt{10.10.10.4/24} и Reception ---
\texttt{10.10.10.5/24}. Фрагмент соединения двух комнат показан на рисунке
14, исходное расположение устройств --- на рисунке 18.

\subsection{Задание 3.6. Изучение межсетевых устройств}

\subsubsection{Часть 1. Физические характеристики}

\paragraph{Шаг 1а. Порты управления.}
Открыта вкладка \textit{Physical} маршрутизатора East, проверено включение
устройства и наличие порта управления.

\paragraph{Шаг 1б. Масштаб и окно устройства.}
Окно Physical увеличено, чтобы увидеть корпус маршрутизатора, расположение
интерфейсов и слотов.

\paragraph{Шаг 1в. Порт управления.}
Определён консольный порт, который используется для локального CLI-доступа
и первоначальной настройки устройства.

\paragraph{Шаг 2а. Интерфейсы LAN и WAN.}
Для East определены физические интерфейсы локальной сети Ethernet и
последовательный WAN-интерфейс. Количество интерфейсов проверено в CLI.

\paragraph{Шаг 2б. Команда show ip interface brief.}
Выполнена команда \texttt{show ip interface brief}. По выводу различены
физические и виртуальные интерфейсы; виртуальный интерфейс существует в
программной конфигурации.

\paragraph{Шаг 2в. Пропускная способность.}
Командой \texttt{show interface gigabitethernet 0/0} определена скорость
интерфейса GigabitEthernet, а командой \texttt{show interface serial 0/0/0}
проверена скорость последовательного интерфейса. Для Serial отображаемая
пропускная способность является параметром маршрутизации, а не обязательно
фактической скоростью линии.

\paragraph{Шаг 3а. Слоты East.}
Определено количество доступных слотов расширения маршрутизатора East.

\paragraph{Шаг 3б. Слоты коммутаторов.}
Проверены Switch2 и Switch3 и определено количество доступных слотов
расширения. Результаты вопросов показаны на рисунке 15.

\subsubsection{Часть 2. Модули и подключение}

\paragraph{Шаг 1а. Модуль для трёх ПК.}
В списке модулей East выбран вариант, позволяющий подключить три
компьютера к маршрутизатору при недостатке встроенных портов. Количество
портов проверено по описанию выбранного модуля.

\paragraph{Шаг 1б. Модуль для Switch2.}
Для Switch2 выбран модуль, обеспечивающий оптоволоконное Gigabit-соединение
с Switch3.

\paragraph{Шаг 2а. Подготовка установки.}
Выбран маршрутизатор East и открыт список доступных слотов.

\paragraph{Шаг 2б. Отключение питания.}
При попытке установки проверено сообщение о невозможности добавить модуль
при включённом питании. East выключен физическим переключателем.

\paragraph{Шаг 2в. Установка и включение.}
Модуль установлен в выбранный слот, после чего питание East включено.
Аналогично подготовлены свободные слоты Switch2 и Switch3.

\paragraph{Шаг 2г. Проверка слота.}
Командой \texttt{show ip interface brief} проверено, в каком слоте
появился новый интерфейс.

\paragraph{Шаг 2д. Маршрутизатор West.}
Для West установлен модуль, добавляющий последовательный интерфейс рядом
с высокоскоростной интерфейсной платой \textit{eHWIC 0}. После установки
проверено появление новых Serial-интерфейсов.

\paragraph{Шаг 3. Подключение устройств.}
Выбраны кабели и интерфейсы согласно таблице соединений. Процедура выбора
первого и второго конца кабеля показана на рисунках 16 и 17.

\subsubsection{Часть 3. Подключение устройств}

Маршрутизатор East соединён с Switch1 через GigabitEthernet, коммутаторы и
компьютеры соединены медными кабелями, а маршрутизаторы East и West —
последовательным DCE-каналом. После подключения проверены индикаторы портов
и соответствие каждой пары интерфейсов таблице на рисунке 17. Скриншоты
заданий по физическим характеристикам, модулям и соединениям приведены на
рисунках 15--18.

\section{Ответы на контрольные вопросы}

\begin{enumerate}
\item Семь уровней OSI/ISO: физический, канальный, сетевой, транспортный,
сеансовый, уровень представления и прикладной.
\item Физический уровень передаёт биты; канальный формирует кадры и использует
MAC-адресацию; сетевой выполняет логическую адресацию и маршрутизацию;
транспортный обеспечивает доставку между процессами; сеансовый управляет
сеансом; уровень представления преобразует формат данных; прикладной
предоставляет сетевые службы приложениям.
\item OSI/ISO является семиуровневой эталонной моделью. TCP/IP объединяет
верхние уровни OSI в прикладной уровень, физический и канальный --- в уровень
доступа к сети, а сетевой уровень называет уровнем Интернета.
\item Физическая топология описывает расположение устройств и кабельные
соединения. Логическая топология описывает путь и правила передачи данных
независимо от физической схемы.
\item Основные среды передачи: витая пара, коаксиальный кабель,
оптоволокно и беспроводной радиоканал.
\item Канальный уровень разделён на LLC и MAC. LLC обеспечивает связь с
верхними уровнями, MAC выполняет инкапсуляцию кадров, MAC-адресацию и
управление доступом к среде.
\item Основные режимы Cisco IOS: пользовательский EXEC, привилегированный
EXEC, глобальный режим конфигурации и специализированные режимы
конфигурации интерфейсов и линий.
\end{enumerate}

\subsection{Ответы по ARP}

\begin{enumerate}[label=\arabic*.]
\item Статическая запись ARP удаляется вручную командой \texttt{arp -d}
с указанием IP-адреса либо вместе с остальными записями командой
\texttt{arp -d *}. Она также может быть удалена при изменении конфигурации
или перезапуске системы, в зависимости от реализации.
\item Статические записи применяются, когда нужно обеспечить постоянное
соответствие IP- и MAC-адреса, уменьшить количество ARP-запросов или
предотвратить подмену ответа ARP для критически важного узла.
\item Полностью отключать ограничения времени хранения записей не следует:
изменения в сети и замена оборудования могут сделать запись устаревшей.
Бессрочная устаревшая запись приведёт к отправке кадров не тому устройству.
\end{enumerate}

\section{Заключение}

В ходе лабораторной работы изучены уровни OSI/ISO, особенности физической и
канальной передачи данных, Ethernet и MAC-адресация. Освоены интерфейс Cisco
Packet Tracer, режимы командной строки Cisco IOS и базовые команды проверки
состояния устройства. Выполнена первоначальная настройка безопасности
коммутатора, рассмотрено сохранение конфигурации в NVRAM. На практике
исследованы формирование ARP-кэша, последовательность ARP-запроса и
ARP-ответа, а также влияние ARP на задержку первого обмена. Полученные навыки
могут использоваться при построении, настройке и диагностике локальных сетей.

\newpage
\section{Иллюстрации выполнения заданий}

\subsection{Знакомство с Packet Tracer}

\paragraph{Рисунок 1.}
На первом скриншоте открыто задание Packet Tracer «Представление сети».
В верхней части окна отображается название практического упражнения, ниже
приведены исходные данные и последовательность действий. В задании
предлагается ознакомиться с рабочей областью Packet Tracer, панелями
устройств и режимами работы. Отдельно указано, что в режиме реального времени
можно строить топологию, а в режиме моделирования --- наблюдать прохождение
пакетов и заголовков протоколов.

\labimage{image1.png}{Интерфейс задания Packet Tracer «Представление сети».}

\paragraph{Рисунок 2.}
На втором скриншоте показана продолженная часть того же упражнения. Здесь
рассматриваются категории оконечных и промежуточных устройств, сетевые
соединения, локальные сети, WAN и Интернет. Вопросы задания требуют
определить назначение устройств, сравнить типы сетей и объяснить различия
между глобальными и локальными соединениями. В нижней части экрана приведены
дополнительные задания по созданию собственной топологии и подключению
оконечных устройств.

\labimage{image2.png}{Изучение локальных сетей, WAN и Интернета.}

\paragraph{Рисунок 3.}
На третьем скриншоте отображена таблица оценивания первого упражнения.
В ней перечислены разделы «Обзор программы Packet Tracer» и «Вселенная
сетей, WAN и Интернет», указано количество вопросов и максимальное число
баллов. Таким образом, этот экран использовался для контроля состава
выполненной работы и объёма проверяемых ответов, а не для настройки отдельного
сетевого устройства.

\labimage{image3.png}{Результаты выполнения задания по представлению сети.}

\subsection{Навигация по Cisco IOS}

\paragraph{Рисунок 4.}
На четвёртом скриншоте открыто упражнение «Навигация по IOS». В первой части
предлагается установить консольное подключение к коммутатору, выбрать
правильный тип кабеля и открыть терминал. Далее требуется перейти в
привилегированный режим и использовать справочную систему Cisco IOS. На
экране также перечислены режимы пользовательского EXEC, привилегированного
EXEC и конфигурационные режимы.

\labimage{image4.png}{Задание Packet Tracer по навигации в Cisco IOS.}

\paragraph{Рисунок 5.}
На пятом скриншоте показано упражнение по режимам ввода команд. В нём
последовательно проверяются команда \texttt{enable}, автодополнение по клавише
Tab, переход в глобальный режим командой \texttt{configure terminal} и
возврат из режима конфигурации командой \texttt{exit}. По изменению приглашения
командной строки можно определить текущий режим: \texttt{Switch>},
\texttt{Switch\#} или \texttt{Switch(config)\#}.

\labimage{image5.png}{Изучение режимов ввода команд Cisco IOS.}

\paragraph{Рисунок 6.}
На шестом скриншоте рассматривается настройка системных часов. Сначала
предлагается вызвать справку по команде \texttt{clock} и определить
обязательные параметры. Затем с помощью знака вопроса уточняется формат даты
и времени, после чего вводится команда \texttt{clock set}. Вторая часть
упражнения показывает дополнительные сообщения IOS и демонстрирует, как
система сообщает об ошибке синтаксиса при неполной команде.

\labimage{image6.png}{Настройка часов и использование справки IOS.}

\paragraph{Рисунок 7.}
На седьмом скриншоте показана итоговая таблица оценивания упражнения по
навигации в IOS. В ней отдельно учитываются базовые подключения, изучение
режимов ввода и настройка часов. Экран позволяет сопоставить выполненные
шаги с максимальным количеством баллов и проверить, какие разделы задания
остались без результата.

\labimage{image7.png}{Результаты выполнения задания по режимам ввода.}

\subsection{Первоначальная настройка коммутатора}

\paragraph{Рисунок 8.}
На восьмом скриншоте открыто упражнение по настройке исходных параметров
коммутатора. В первой части предлагается перейти в привилегированный режим,
выполнить \texttt{show running-config} и изучить конфигурацию устройства по
умолчанию. Отдельно проверяются количество FastEthernet- и GigabitEthernet-
интерфейсов, значения параметров по умолчанию и отсутствие файла
\texttt{startup-config} до первого сохранения.

\labimage{image8.png}{Проверка конфигурации коммутатора по умолчанию.}

\paragraph{Рисунок 9.}
На девятом скриншоте показаны команды настройки имени коммутатора и защиты
консольного доступа. В глобальном режиме задаётся имя устройства, затем
выполняется переход в режим \texttt{line console 0}, задаётся пароль и
включается команда \texttt{login}. После выхода из режима конфигурации
проверяется появление сообщения о необходимости ввести пароль при новом
подключении к консоли.

\labimage{image9.png}{Настройка основных параметров и доступа к коммутатору.}

\paragraph{Рисунок 10.}
На десятом скриншоте продолжается настройка безопасности. Задаётся пароль
привилегированного режима с помощью \texttt{enable password}, затем он
проверяется повторным входом. После этого используется более безопасная
команда \texttt{enable secret}, включается шифрование паролей командой
\texttt{service password-encryption} и настраивается предупреждающее MOTD-
сообщение. Вопросы внизу экрана требуют объяснить назначение команды
\texttt{login} и необходимость баннера для всех пользователей устройства.

\labimage{image10.png}{Проверка паролей, шифрование и настройка MOTD.}

\paragraph{Рисунок 11.}
На одиннадцатом скриншоте показано сохранение текущей конфигурации в NVRAM.
Сначала просматривается текущая конфигурация командой \texttt{show run},
после чего выполняется \texttt{copy running-config startup-config}. Затем
проверяется содержимое загрузочной конфигурации и выполняется аналогичная
настройка второго коммутатора: имя, пароль консоли, пароль привилегированного
режима, MOTD и сохранение параметров.

\labimage{image11.png}{Сохранение конфигурации и настройка второго коммутатора.}

\paragraph{Рисунок 12.}
На двенадцатом скриншоте приведена таблица оценивания задания по настройке
коммутатора. Отдельными строками проверяются просмотр конфигурации по
умолчанию, имя устройства, защита консоли, пароли привилегированного режима,
шифрование, MOTD, сохранение в NVRAM и настройка второго коммутатора. Внизу
экрана отображаются текущие время выполнения, набранные баллы и максимальная
оценка; эти значения позволяют отличить выполненные шаги от ещё не
проверенных.

\labimage{image12.png}{Результаты выполнения задания по настройке коммутатора.}

\subsection{Комплексная настройка сети}

\paragraph{Рисунок 13.}
На тринадцатом скриншоте показано комплексное упражнение с таблицей
адресации. Для устройств Room-145 и Room-146 указаны VLAN-интерфейсы и
адреса \texttt{10.10.10.100/24} и \texttt{10.10.10.150/24}. Для конечных
устройств Manager и Reception приведены адреса \texttt{10.10.10.4/24} и
\texttt{10.10.10.5/24}. Ниже перечислены требования: задать имена
устройств, настроить адресацию, пароли и MOTD, сохранить конфигурации и
проверить наличие соединения между всеми устройствами.

\labimage{image13.png}{Комплексная настройка: таблица адресации и требования.}

\paragraph{Рисунок 14.}
На четырнадцатом скриншоте показан фрагмент построенной топологии. Компьютер
Manager подключён к коммутатору Room-145, а компьютер Reception --- к
коммутатору Room-146. Между коммутаторами отображается межкоммутаторное
соединение. Зелёные индикаторы портов показывают активное состояние
соединений, а цветные линии позволяют различить выбранные и физические
соединения в рабочей области Packet Tracer.

\labimage{image14.png}{Фрагмент топологии Room-145 и Room-146.}

\subsection{Изучение межсетевых устройств}

\paragraph{Рисунок 15.}
На пятнадцатом скриншоте открыто задание по изучению межсетевых устройств.
В первой части требуется исследовать физические характеристики маршрутизатора
East: наличие питания, объём памяти и доступные интерфейсы. Затем через CLI
выполняются команды \texttt{show ip interface brief},
\texttt{show interface gigabitEthernet 0/0} и
\texttt{show interface serial 0/0/0}. По результатам необходимо определить
количество физических интерфейсов, пропускную способность и наличие слотов
расширения.

\labimage{image15.png}{Определение физических характеристик межсетевых устройств.}

\paragraph{Рисунок 16.}
На шестнадцатом скриншоте показан выбор модулей расширения и подключение
оборудования. Для маршрутизаторов подбираются модули, добавляющие
необходимые последовательные или GigabitEthernet-интерфейсы. Перед установкой
модуля устройство требуется выключить, а после установки снова включить.
Далее выбираются тип кабеля и конкретные порты для соединения маршрутизатора,
коммутатора и конечных устройств.

\labimage{image16.png}{Выбор модулей и подключение устройств.}

\paragraph{Рисунок 17.}
На семнадцатом скриншоте приведена таблица требуемых соединений. В каждой
строке указаны устройство и интерфейс с одной стороны, тип кабеля, а также
устройство и интерфейс с другой стороны. В таблице встречаются прямые
медные кабели для Ethernet-соединений, перекрёстный медный кабель для
соединения однотипных устройств и последовательный кабель DCE для
межмаршрутизаторного канала. Эта таблица используется как контрольная схема
при физическом подключении всех компонентов.

\labimage{image17.png}{Таблица соединений межсетевых устройств.}

\paragraph{Рисунок 18.}
На восемнадцатом скриншоте показана исходная рабочая область комплексного
задания до соединения устройств. Видны маршрутизаторы West и East,
коммутаторы Switch1--Switch4 и конечные компьютеры PC1--PC9. Расположение
устройств по группам соответствует дальнейшей таблице соединений: сначала
подготавливаются маршрутизаторы и коммутаторы, затем к ним подключаются
конечные узлы и проверяется связность.

\labimage{image18.png}{Схема расположения устройств в комплексном задании.}

\end{document}
